\section{Development cases and input constructions}
\label{sec:development-details}
This appendix gives the inputs, source requirements, and repair observations
for the three design studies in Sections~\ref{sec:evidence}--\ref{sec:feedback-results}.
\subsection{Requirement memory: preserving triggering conditions}
\label{sec:details-evidence}
\label{sec:details-content}
The executable-memory experiments use a Verified split of 280 source tasks,
70 development tasks, and 150 held-out evaluation tasks. The cases below are
development interventions with DeepSeek V4 Pro through OpenRouter,
temperature zero and matched budgets within each comparison. They measure
whether history adds a required behavior to the submitted repair.
Table~\ref{tab:historicalrequirements} presents three such additions.
Django14404 uses the first-ranked source issue and patch; Django12708 and
Django11138 come from a fixed six-task cohort receiving three unchanged Flat
rules per task. Both arms pass official verification on all six cohort tasks.
The historical requirements distinguish the completeness of those repairs.

\paragraph{An omitted behavior of the same API.}
Django14404 reports that administrative redirects lose a deployment prefix.
Both agents fix the prefix and pass official verification, but only the repair informed by source16612 also preserves the query string.
A joint input combines the prefix with repeated and percent-encoded query values, \texttt{keep=a\%2Fb}\allowbreak\texttt{\&keep=c\%20d}.
Only the history-informed repair preserves both.

\paragraph{A neighboring operation.}
Django12708 reports ambiguous deletion of an ordinary index when a unique constraint uses the same columns.
The first retrieved source, 15732, instead concerns deletion of a unique constraint sharing its column with a primary key.
Both repairs fix ordinary-index selection; only the memory-informed repair also excludes primary keys when deleting unique constraints.
The source-generated PostgreSQL check runs unchanged on both target repairs.
A separate combined input creates an ordinary index, redundant uniqueness, and a primary key on the same column.
Both repairs preserve ordinary-index deletion, but only the memory-informed repair also deletes redundant uniqueness while retaining the primary key and data.
History supplies a requirement for the neighboring operation.

\paragraph{An independent argument of the same operation.}
Django11138 requires date conversions to respect the database's source timezone.
Source13363 requires \texttt{TruncDate} and \texttt{TruncTime} to honor an expression's explicit destination timezone.
Both agents repair the source timezone, but only the memory-informed repair also honors the destination.
A source-generated check distinguishes the historical versions and runs unchanged on the target repairs.
On an independent SQLite input combining a UTC+1 database with an explicit New York destination, only the memory-informed repair returns the correct date and time.
The issue and history constrain different ends of the conversion.

\paragraph{A requirement can suffice without the source implementation.}
On the redirect task, Flat memory and its original issue/patch replacement
produce identical repairs preserving the joint behavior. Two new agents then
revise the ordinary patch with a source-generated requirement: the redirect
must include the original query string. One also receives the check and its
failed observations. Both submit identical repairs passing official and joint
checks. The sentence-only agent finds \texttt{get\_full\_path()} in the current
repository: memory supplies the behavior, and the agent locates its implementation.

\paragraph{Which information makes the requirement usable?}
We select the constraint and timezone benefit cases and run four fresh complete agents per task, fixing the single historical source and its delivery before editing.
Two source-only constructions produce a behavioral requirement and an implementation reference; the combined input concatenates them unchanged.
All eight repairs pass official verification (Table~\ref{tab:memorycontent}).
On Django12708, only Implementation and Both also remove redundant uniqueness while preserving the primary key.
The Requirement agent tests a field with \texttt{unique=True}, but places \texttt{unique\_together} on a different, two-column set.
This nearby passing test loses the same-column condition stated in memory.
On Django11138, all three memory inputs preserve the explicit destination timezone, while the fresh no-memory repair omits it.
The Requirement agent finds \texttt{self.get\_tzname()} in the target implementation without being given that helper in memory.



The timezone repairs also differ on further conditions.
Both still generates an unnecessary MySQL conversion between identical zones; Implementation instead misinterprets a named Paris source timezone, producing New York time \texttt{19:21} instead of \texttt{17:30} on a summer input.
Requirement passes both additional diagnostics.
These concrete omissions distinguish carrying a requirement, applying its triggering condition, and implementing the resulting behavior correctly.
Two subsequent original-base agents receive the unchanged constraint Requirement, with or without a source-validated example; both pass official and joint checks.
The same requirement succeeds when its trigger is reconstructed correctly; preserving that input construction is a useful role for the graph.

\paragraph{Preserve the input relation.}
On Django12663, both a three-source code-connected packet and a ten-source
similarity packet contain the historical named-tuple repair, source13590,
yet neither initial repair recovers its behavior. Three fresh revisions start
from the same candidate and source (Figure~\ref{fig:memory-content-scope}).
Original history and execution feedback recover named tuples; the requirement
sentence alone does not. The feedback repair expands every non-builtin
sequence into positional arguments, which also breaks ordinary tuple and list
subclasses. Original history retains the constructor distinction and passes
all displayed behaviors. This is precisely the relation the graph should carry:
the argument convention of the input's type, alongside its endpoint values.

\begin{figure}[t]
\centering
\includegraphics[width=\linewidth]{figures/memory_content_scope.pdf}
\caption{Memory content changes which behaviors a repair preserves. One revision per input starts from the same Django12663 candidate and source13590. Parentheses give case counts; each cell reports all cases in its group. The ordinary-subclass inputs were chosen after inspecting the feedback repair. All three repairs pass official verification, while only original history passes every displayed behavior.}
\label{fig:memory-content-scope}
\end{figure}

\paragraph{Retain candidates that can supply an omitted condition.}
In this same case, demanding an operation match from the initial public-code
localization excludes source13590, which FAISS ranks first. Keeping the
similarity candidate lets the graph return its input relation and construct
the joint program described next. This motivates the division in Figure~\ref{fig:method}:
FAISS proposes sources; graph relations support their adaptation to the current
operation. The usefulness of a candidate is established by what its condition
adds to the repair.

\subsection{Condition composition: checking interactions}
\label{sec:details-composition-results}

\paragraph{Join the historical input to the current operation.}
The current issue involves a lazy User in a nested query; history contributes
a named-tuple range. On a separate history-informed repair, the current scalar query and
the historical range query each pass, while a query combining them fails.
The composed condition reveals an omission in the second container
reconstruction.

The deterministic adapter in Equation~\ref{eq:composition} reaches the same
omission while retaining the public program's lazy User and nested annotation.
It produces a named-tuple point range and an ordinary-tuple neighbor.
Graph delivery subsequently reconstructs this same program from the retrieved
input relation. It operationalizes the condition tested in the repair comparison.

\paragraph{A matched comparison tests the added condition and feedback.}
A new original-base pair receives identical history, public scalar program,
its observed failure, a callable execution tool, and instructions to run the
supplied conditions before editing and submission. The augmented arm also
receives the two composed programs and their failures.
Both repairs pass official verification (Table~\ref{tab:composedconditions}).
The current-feedback repair fails both historical examples and the withheld
named-tuple/lazy-User query; the augmented repair passes those behaviors and
preserves the ordinary subclasses.
The augmented agent runs the joint input before editing and again after
removing a global conversion that broke an existing test. Its remaining
range-endpoint failure directs the final correction.

\paragraph{Compose roles to expose a broader requirement.}
The timezone relation provides another example. A fixed adapter varies the
historical witness's three zones across database, ambient, and destination
roles and shifts the input by 182 days. Of 72 combinations, the source patch
changes 20 from fail to pass; 20 pass on both versions and 32 fail on both.
Applying the family to the four timezone content interventions shows that Requirement and
Both pass all 72. Implementation passes the 40 source-supported combinations
but fails the remaining 32, including 16 passed without memory.
Role composition exposes how the current source-timezone requirement interacts
with the historical destination requirement. A complete repair can therefore
satisfy combinations beyond those corrected by the historical implementation.

\subsection{Execution feedback: guiding requirement-based repair}
\label{sec:details-feedback}

\paragraph{The addition can guide further repair.}
Two fresh agents start from the same completed ordinary timezone patch.
Issue-based additional review returns it unchanged; supplying a generated historical check and its failed observations leads the other agent to use \texttt{self.get\_tzname()} in both truncation methods.
Both revisions pass official verification, but only the latter passes the historical and joint inputs.
The intervention supplies corrective information that further issue-based review does not recover in this comparison.

\paragraph{Execution feedback makes an omitted condition concrete.}
Two fresh agents refine the same original failed Requirement candidate, both receiving its requirement again.
One also receives the unchanged source-generated check and its actual failures for a field with a primary key or \texttt{unique=True}.
Requirement-only review excludes the primary key but still fails the ordinary unique-field condition.
Feedback adds generated-name disambiguation and passes both conditions.
Both revisions pass official verification and the withheld joint sequence.
The added behavior is recovery of the ordinary unique-field condition identified by feedback.

\paragraph{Feedback from the composed query.}
The Django12663 example also provides a refinement comparison.
Four fresh refinements of its history-informed candidate illustrate the use of feedback:
history alone yields one of two repairs passing all seven measured conditions;
adding the constructed joint program and its failed observation yields two of
two. All four pass official verification. Both augmented agents execute the
joint input and repair the second reconstruction.
